% MAS1004 Assignment 2: long report.
%
% You fill this in with your agent, in your repository. AGENTS.md tells the
% agent how. Red text in the PDF is what still has to be replaced.
%
% - Pictures go in figures/. Copy them from results/, which is not in git.
% - Outputs go in outputs/: copy experiments.csv, LOG.md and the outputs of
%   check_submission.py and curate.py count here. \outputfile{outputs/...}
%   prints a file exactly as it is. Every number in this report has to come
%   from such an output.
%
% To make the PDF: compress this folder into a zip, and on Overleaf choose
% New Project, Upload Project. Set Menu, Compiler to XeLaTeX, and Recompile.

\documentclass[11pt]{article}
\usepackage{mas1004}

\reporttitle{Assignment 2: Long report}
\studentname{Kim Sofia}
\studentid{20251232}

\begin{document}
\makeheader

Leaderboard nickname: 허니머스타드. Current releases: Track A \texttt{track-a-1} (starter recipe) and Track B \texttt{track-b-1} (Track B starter baseline). Track A leaderboard score: Not observed. Track B leaderboard score: Not observed.

% The sections below are a suggestion. Change, merge or add sections as you
% like; README.md, "The long report", says what a reader should be able to see.

% ---------------------------------------------------------------------------
\section*{Track A}

\begin{center}\footnotesize
\begin{tabular}{lp{0.6\linewidth}r}
\toprule
tag & what changed from the run before & validation \\
\midrule
baseline, seed 0 & Starter recipe; all 44,257 starter training photos; official budget of 249,984 photo exposures. & 0.4967 \\
baseline, seed 1 & Same starter recipe and data as seed 0; seed changed only to measure run-to-run variation. & 0.4948 \\
\bottomrule
\end{tabular}
\end{center}

The two baseline scores differed by 0.0019 (mean 0.4958). I used these two runs as a starting measurement and as an estimate of seed-to-seed variation. No cleaning, selection, augmentation, or architecture conclusion is drawn from the baseline alone.

The current Track A release is \texttt{track-a-1}, the unchanged starter recipe. Its leaderboard score was not observed, so no comparison between leaderboard and validation is claimed.

% ---------------------------------------------------------------------------
\section*{Decisions you made from photos}

I did not copy class-only starter-photo contact sheets into this public repository. I inspected representative training sheets and the locally generated b1 mistake and suspect sheets for Busan and Osaka.

In the suspicious training images, I observed some isolated flowers, birds, food, museum objects, and indoor spaces without visible urban surroundings. In the error sheets, Osaka street and traditional-building scenes were often predicted as Tokyo, while Busan scenes were often predicted as Seoul. The saved Track B b1 confusion matrix recorded 785 Osaka validation photos predicted as Tokyo and 410 Busan validation photos predicted as Seoul.

I wrote the following prospective rule before acting: exclude only training photographs that show isolated objects, food, animals, or other close-up subjects without recognizable urban or location-specific context. This rule is \textbf{untested}: no photo was removed, no selection file was created, and no follow-up model was trained before the deadline.

The relevant requests were: ``Some suspicious training images contain almost no useful information about their city'' and ``Exclude only training photographs that show isolated objects, food, animals, or other close-up subjects without any recognizable urban or location-specific context.'' I also specified that the model architecture, augmentation, random seed, and official training budget should stay unchanged for a future comparison. No action was taken on that proposal.

\subsection*{Saved observation and decision record}
\outputfile{outputs/prompts_and_decision.txt}

% ---------------------------------------------------------------------------
\section*{Track B}

\begin{center}\footnotesize
\begin{tabular}{lp{0.6\linewidth}r}
\toprule
tag & what changed from the run before & validation \\
\midrule
b1 & Starter Track B baseline: ResNet18 at 224 pixels, ImageNet-1K V1 starting weights, \texttt{data/starter/train}, 4 epochs. & 0.4961 \\
\bottomrule
\end{tabular}
\end{center}

The b1 submission passed the local format and limit checker: 11.17 million parameters, 1.81 billion multiply-adds per 224-by-224 photo, and a 44.7 MB ONNX file. This established a valid starter Track B submission before attempting a controlled improvement. No changed Track B model was trained.

No extra photos were used in b1; it trained on \texttt{data/starter/train}.

The b1 validation score was 0.4961. Its leaderboard score was not observed, so no comparison between leaderboard and validation is claimed.

\subsection*{Your submission, as the grader sees it}
\outputfile{outputs/check_submission.txt}

% ---------------------------------------------------------------------------
\section*{What made the difference}

No changed cleaning, augmentation, or additional-data experiment was completed before the deadline. Therefore, this report does not claim that any change helped or hurt.

With one more week, I would test the written close-up/isolated-subject exclusion rule as one controlled Track A change, keeping the architecture, augmentation, seed comparison, and official budget fixed. I would then compare its validation score with the baseline and inspect the new errors before making another change.

% ---------------------------------------------------------------------------
\section*{All runs}
\outputfile{outputs/experiments.csv}

\section*{Lab notebook}
\outputfile{outputs/LOG.md}

\end{document}
